\section{Introducción}
\label{sec:intro}

\subsection{Contexto}

La comunicación digital se ha convertido en el canal primario de coordinación en
entornos educativos. Grupos de mensajería instantánea como Telegram, WhatsApp o Discord
albergan hoy una parte sustancial de la interacción informal entre estudiantes,
docentes y personal administrativo. Estos espacios no solo transmiten información
logística ---fechas, requisitos, procedimientos--- sino que también vehiculan estados
afectivos: dudas ansiosas ante plazos, gratitud por la ayuda recibida, frustración ante
trámites fallidos o entusiasmo colectivo por logros compartidos.

El proyecto AS\_Telegram se propone analizar, de manera exploratoria y sistemática, la
carga emocional y la estructura temática de uno de esos espacios: el grupo privado de
Telegram ``Curso de Temporada Invierno 2026'', dedicado a la gestión académica de un
programa universitario. El objetivo no es evaluar a personas ni instituciones, sino
caracterizar el clima emocional agregado y comprender qué tipos de intercambio dominan
la conversación.

\subsection{Problema de investigación}

El análisis de sentimiento en español aplicado a contextos operativos plantea tres
dificultades que motivan este trabajo. En primer lugar, la mayor parte del contenido de
un grupo de gestión académica es \emph{factual} y no \emph{afectivo}: preguntas por
horarios, enlaces a formularios o confirmaciones de pago rara vez contienen palabras
cargadas de sentimiento. Un método que fuerce una etiqueta en todos los mensajes
producirá una proporción artificialmente alta de \emph{neutro}, confundiendo la ausencia
de señal con la neutralidad. En segundo lugar, el texto informal está poblado de emojis,
jerga, abreviaturas y errores tipográficos que los métodos entrenados sobre texto formal
manejan mal. En tercer lugar, la negación y las construcciones irónicas pueden invertir
el sentido de un término léxico.

Este informe sostiene que un análisis honesto debe (i) distinguir explícitamente entre
\emph{ausencia de evidencia} y \emph{neutralidad}, (ii) aprovechar los emojis como
señal informativa en lugar de descartarlos y (iii) reportar la cobertura del método como
un resultado de primer orden, no como una nota al pie.

\subsection{Objetivos}

El objetivo general es dise\~nar, ejecutar y documentar un \emph{pipeline} reproducible
de análisis de sentimiento y emociones para un corpus real de Telegram en español. Los
objetivos específicos son:

\begin{enumerate}[label=\textbf{O\arabic*.}]
  \item Construir una infraestructura ETL trazable que importe, valide, limpie y
        persista el corpus, registrando cada operación en una base de datos de
        auditoría.
  \item Desarrollar un clasificador léxico multietiqueta en español que integre las
        señales del léxico, los emojis y las reacciones, y que modele la negación de
        forma contextual.
  \item Caracterizar la distribución de polaridad y emociones, prestando especial
        atención a la cobertura del diccionario y a los mensajes indeterminados.
  \item Describir la evolución temporal y los patrones de actividad (día y hora).
  \item Descubrir la estructura temática latente del corpus mediante LDA.
  \item Modelar las transiciones entre estados de polaridad mediante cadenas de Markov.
  \item Generar de forma automática un informe reproducible que se actualice con los
        datos.
\end{enumerate}

\subsection{Preguntas de investigación}

El estudio se guía por las siguientes preguntas:

\begin{description}[style=nextline,leftmargin=1.2cm]
  \item[PI1.] ¿Qué proporción del corpus contiene evidencia afectiva suficiente y, por
        tanto, qué fracción queda indeterminada?
  \item[PI2.] Entre los mensajes con evidencia, ¿cuál es el balance entre polaridad
        positiva y negativa, y qué emociones predominan?
  \item[PI3.] ¿Cómo se distribuye la actividad y la polaridad a lo largo del tiempo?
  \item[PI4.] ¿Qué temas latentes estructuran la conversación?
  \item[PI5.] ¿Cómo se transita entre estados emocionales y qué estados son más
        persistentes?
\end{description}

\subsection{Contribuciones}

Este trabajo realiza cuatro contribuciones. Primero, ofrece un \emph{pipeline} completo,
abierto y reproducible para el análisis de sentimiento en grupos de Telegram en español,
incluyendo el tratamiento de emojis, reacciones y negación. Segundo, propone una
convención explícita para el manejo de la evidencia que evita la imputación de
neutralidad. Tercero, aplica el método a un corpus real de más de veinte mil mensajes,
ofreciendo estadísticas descriptivas detalladas. Cuarto, integra la generación
automática del informe en el propio \emph{pipeline}, de modo que las cifras del documento
se derivan directamente de los datos.

\subsection{Estructura del documento}

El resto del informe se organiza así. La Sección~\ref{sec:marco} presenta el marco
teórico. La Sección~\ref{sec:metodologia} describe la metodología. La
Sección~\ref{sec:datos} caracteriza los datos. La Sección~\ref{sec:resultados} presenta
los resultados. La Sección~\ref{sec:discusion} los discute. La
Sección~\ref{sec:limitaciones} aborda las limitaciones. La Sección~\ref{sec:conclusiones}
concluye. La Sección~\ref{sec:etica} trata la ética, la
Sección~\ref{sec:reproducibilidad} la reproducibilidad, y el apéndice recoge material
complementario.
